\documentclass[10pt,a4paper]{article}

\usepackage[margin=30mm]{geometry}
\usepackage{amsmath,amssymb,amsthm,mathtools}
\usepackage{enumitem}
\usepackage{array}
\usepackage[pdftex,
pdftitle={Exact finite enumeration of a bounded stereochemical SMILES sublanguage},
pdfauthor={},
colorlinks,
linkcolor=black,
urlcolor=black,
citecolor=black]{hyperref}

\allowdisplaybreaks
\emergencystretch=3em

\newcommand{\calA}{\mathcal A}
\newcommand{\calB}{\mathcal B}
\newcommand{\calC}{\mathcal C}
\newcommand{\calD}{\mathcal D}
\newcommand{\calE}{\mathcal E}
\newcommand{\calF}{\mathcal F}
\newcommand{\calG}{\mathcal G}
\newcommand{\calH}{\mathcal H}
\newcommand{\calI}{\mathcal I}
\newcommand{\calK}{\mathcal K}
\newcommand{\calL}{\mathcal L}
\newcommand{\calM}{\mathcal M}
\newcommand{\calO}{\mathcal O}
\newcommand{\calP}{\mathcal P}
\newcommand{\calQ}{\mathcal Q}
\newcommand{\calR}{\mathcal R}
\newcommand{\calS}{\mathcal S}
\newcommand{\calT}{\mathcal T}
\newcommand{\calV}{\mathcal V}
\newcommand{\calX}{\mathcal X}
\newcommand{\calY}{\mathcal Y}
\newcommand{\calZ}{\mathcal Z}
\newcommand{\Nat}{\mathbb N}
\newcommand{\Bool}{\{0,1\}}
\newcommand{\eps}{\varepsilon}
\newcommand{\dom}{\operatorname{dom}}
\newcommand{\im}{\operatorname{im}}
\newcommand{\Aut}{\operatorname{Aut}}
\newcommand{\Sol}{\operatorname{Sol}}
\newcommand{\render}{\operatorname{render}}
\newcommand{\parse}{\operatorname{parse}}
\newcommand{\supp}{\operatorname{supp}}
\newcommand{\sgn}{\operatorname{sgn}}
\newcommand{\pos}{\operatorname{pos}}
\newcommand{\None}{\bot}
\newcommand{\Ring}{\mathrm{ring}}
\newcommand{\Tree}{\mathrm{tree}}
\newcommand{\br}{\mathrm{br}}
\newcommand{\cont}{\mathrm{cont}}
\newcommand{\dotcup}{\mathbin{\dot\cup}}
\newcommand{\st}{\;:\;}
\newcommand{\Ann}{\mathsf{Ann}}
\newcommand{\OK}{\mathsf{OK}}
\newcommand{\Dir}{\mathsf{Dir}}
\newcommand{\Tet}{\mathsf{Tet}}
\newcommand{\St}{\mathsf{St}}
\newcommand{\Site}{\mathsf{Site}}
\newcommand{\Decode}{\mathsf{Decode}}
\newcommand{\Lin}{\mathsf{Lin}}
\newcommand{\Base}{\mathsf{Base}}
\newcommand{\Max}{\mathsf{Max}}

\newtheorem{lemma}{Lemma}[section]
\newtheorem{proposition}[lemma]{Proposition}
\newtheorem{theorem}[lemma]{Theorem}
\newtheorem{corollary}[lemma]{Corollary}
\theoremstyle{definition}
\newtheorem{definition}[lemma]{Definition}
\newtheorem{assumption}[lemma]{Assumption}
\newtheorem{example}[lemma]{Example}
\theoremstyle{remark}
\newtheorem{remark}[lemma]{Remark}

\title{Exact finite enumeration of a bounded stereochemical SMILES sublanguage}
\author{}
\date{\today}

\begin{document}
	\maketitle

	\begin{abstract}
		We formulate exact support enumeration for a bounded RDKit/OpenSMILES-like
		stereochemical SMILES sublanguage as a finite declarative problem.  The input is
		a fixed finite molecular graph whose atom and bond attributes, aromaticity,
		implicit and explicit hydrogens, formal charges, isotopes, bond orders, and
		specified stereochemistry are already fixed.  The language is generated by an
		attributed graph-traversal grammar, and the remaining choices are finite-domain
		syntax-slot variables: atom text, bond text, ring labels, branch and
		continuation choices, tetrahedral tokens, and directional slash/backslash
		markers.  A finite constraint system enforces graph identity, parser-defined
		stereo preservation, absence of accidental stereo, ring-label correctness, and
		the declared annotation policy.  The exact support is the image of the rendering
		map on satisfying assignments.  A canonical-witness constraint may be added for
		non-redundant enumeration, but it is not needed for the mathematical definition
		of the support.  Under the stated bounded and explicit-semantics assumptions, no
		post-hoc parse filter, writer-specific repair rule, or feature-specific oracle is
		forced.
	\end{abstract}

	\noindent {\it Keywords}: SMILES, stereochemistry, enumeration, constraint
	satisfaction, attributed grammar, finite language, tetrahedral chirality,
	directional bonds, ring closures.

	\section{Question and main answer}
	\label{s:intro}

	Fix a finite molecular graph whose non-stereochemical atom and bond attributes
	are already determined.  Fix also a finite SMILES sublanguage: finite atom-text
	and bond-text choices, bounded normalized ring labels, a fixed fragment-order
	policy, and an explicit policy for tetrahedral and directional-bond annotation.
	The question is whether the exact support of that sublanguage can be enumerated
	as one finite declarative object over traversal events and syntax-slot
	assignments, with validity and stereochemical parse-equivalence guaranteed by
	construction.

	The answer is yes.  The clean formulation is
	\[
	\text{finite attributed graph-traversal grammar}
	\quad + \quad
	\text{finite-domain CSP/SAT constraints} .
	\]
	The grammar generates exactly the allowed traversal-event skeletons.  The finite
	constraints assign text and stereo values to emitted syntax slots and require
	that the declared parser semantics sends the rendered string to the input
	labelled stereochemical graph.  A pure finite automaton also exists after full
	compilation, because the instance language is finite, but the automaton view
	hides the structure.  A plain context-free grammar is not the natural object,
	because ring pairing, graph coverage, local tetrahedral ligand order, shared
	slash/backslash carriers, and maximal annotation policies are finite relational
	constraints rather than context-free syntax alone.

	The conclusion is relative to a declared semantics.  If ``valid'' means
	``accepted after implementation-specific sanitization, writer repair, stereo
	cleanup, aromaticity reperception, or canonicalization'', then those behaviours
	are part of the mathematical target.  They may still be finite in a bounded
	setting, but they must be encoded as relations or explicitly treated as oracles.
	The bounded dialect itself does not force such oracles.

	\section{Standing assumptions}
	\label{s:assumptions}

	Let
	\[
	G=(V,E,\alpha,\beta,\chi,\omega)
	\]
	be a fixed finite labelled molecular graph.  The graph is taken to be undirected
	and simple only for notational economy; molecular multigraph conventions and
	pseudo-ligand occurrences can be represented by replacing the corresponding
	finite index sets.  The maps \(\alpha\) and \(\beta\) contain the fixed
	non-stereochemical atom and bond attributes.  The partial map \(\chi\) specifies
	tetrahedral stereo at a finite set \(T\subseteq V\).  The partial map \(\omega\)
	specifies ordinary directional double-bond stereo at a finite set of stereo
	sites, usually represented by eligible double bonds in \(E\).  Enhanced stereo
	groups, non-tetrahedral atom stereo, atropisomerism, unknown stereo,
	unspecified stereo as an input annotation, and unbounded ring-label syntax are
	outside the scope.

	The fixed policy \(P\) consists of finite data:
	\begin{enumerate}[label=\textup{(P\arabic*)},leftmargin=2.5em]
		\item a finite traversal-event language: roots, rooted spanning forests, branch
		order, continuation-child choice, ring-closure endpoint positions, and a fixed
		ordering of disconnected fragments;
		\item for each atom occurrence, a finite set of permitted atom texts;
		\item for each emitted bond or ring endpoint occurrence, a finite set of
		permitted bond texts and a finite convention for possible directional markers;
		\item a finite ring-label alphabet \(\Lambda\), together with the text used for
		each label;
		\item finite parser relations for the bounded fragment: non-stereo atom and bond
		decode, tetrahedral \(@\)/\(@@\) decode, directional slash/backslash decode, and
		ring-label pairing;
		\item an annotation policy, for example support-wise maximal eligible
		directional carrier annotation, cardinality-maximal annotation, a canonicalized
		variant, or a hard predeclared set of required annotation slots.
	\end{enumerate}

	The parser semantics is assumed to be a finite partial function
	\[
	\parse_P:\Sigma_P^*\rightharpoonup \mathcal G_P
	\]
	from legal token strings in the bounded dialect to labelled stereochemical
	molecular graphs.  Equivalently, one may use a finite parse relation rather than
	a function; then the theorem below is read with the chosen existential or
	universal notion of parse-equivalence.  The functional presentation is cleaner
	and avoids hiding nondeterminism in the proof.  Ring closures, implicit
	hydrogens, aromaticity, and directional-bond interpretation are therefore not
	implicit implementation calls: they are either fixed in \(G\), fixed in \(P\), or
	represented by finite relations in \(P\).

	\begin{remark}[Bounded text choices]
		The finiteness assumption is stronger than merely bounding ring labels.  Any
		integer-like or extensible text field that is admitted into the language must be
		bounded or fixed: isotope spelling, atom-map numbers, atom classes, CXSMILES
		fields, aliases, charge spellings if the grammar permits redundant variants,
		and any implementation extension that carries arbitrary text.  Once each atom,
		bond, ring-label, and extension field has a finite domain, allowing all
		syntactically distinct but semantically equivalent spellings does not threaten
		finiteness.
	\end{remark}

	\section{Traversal skeletons}
	\label{s:skeleton}

	The syntactic backbone is most cleanly described as an ordered depth-first
	traversal skeleton.  Equivalently, it is a rooted spanning forest together with
	an order of child and cotree-edge endpoint events at each vertex.

	Let \(V_1,\ldots,V_m\) be the connected components of \(G\), in the fragment
	order fixed by \(P\).  A traversal skeleton \(S\) consists of:
	\begin{enumerate}[label=\textup{(S\arabic*)},leftmargin=2.5em]
		\item a root \(r_i\in V_i\) for each component;
		\item a parent map
		\[
		p:V\to V\cup\{\None\}
		\]
		with \(p(r_i)=\None\), with \(p(v)\) adjacent to \(v\) for each non-root vertex,
		and with the directed parent relation acyclic and reaching the component root;
		\item the induced partition
		\[
		E=E_{\Tree}\dotcup E_{\Ring},
		\qquad
		E_{\Tree}:=\{\{v,p(v)\}:p(v)\ne\None\};
		\]
		\item for each vertex \(v\), a finite ordered list \(O_v\).  Its entries are the
		ring endpoint items \((e,v)\) for cotree edges \(e\in E_{\Ring}\) incident with
		\(v\), and the child items \(u\) with \(p(u)=v\).  At most one child item is
		marked as the continuation child; every other child item is a branch child.
	\end{enumerate}
	The continuation marker is a presentation choice, not a graph-theoretic edge
	property.  One can include it explicitly in the skeleton, as above, or encode it
	as part of a linearized DFS event word with push and pop events.  These two
	views are equivalent: the event word determines which recursive descent is
	unparenthesized, and the marked skeleton determines the event word.  Keeping the
	marker explicit makes the later constraints easier to read.

	Given a skeleton and slot values, the recursive renderer has the form
	\[
	R(v)=\operatorname{atom}_v\;\mathop{\Vert}_{o\in O_v}\Phi(o),
	\]
	where \(\Vert\) denotes concatenation in the order of \(O_v\), and
	\[
	\Phi(o)=
	\begin{cases}
		\operatorname{bondring}_{e,v}\,\lambda_e,
		& o=(e,v)\text{ is a ring endpoint},\\[0.3em]
		(\operatorname{bond}_{v,u}\,R(u)),
		& o=u\text{ is a branch child},\\[0.3em]
		\operatorname{bond}_{v,u}\,R(u),
		& o=u\text{ is the continuation child}.
	\end{cases}
	\]
	The displayed formula is schematic.  The atom and bond expressions are filled by
	finite slot variables.  The full string is
	\[
	R(r_1).R(r_2).\cdots.R(r_m),
	\]
	with dots omitted when \(m=1\).  The set of skeletons satisfying these
	conditions and any additional traversal normalizations in \(P\) is denoted by
	\(\calS(G,P)\).  It is finite.

	\begin{remark}[Canonical graph-theoretic object]
		The minimal graph-theoretic object is not merely a rooted spanning forest.  The
		same forest can yield different SMILES strings depending on the order of ring
		endpoint events, the order of branch children, and the choice of continuation
		child.  A more canonical object is therefore an ordered DFS event skeleton: a
		rooted spanning forest plus a total local order of all child and cotree-endpoint
		events, with enough information to recover the balanced-parenthesis
		linearization.
	\end{remark}

	\section{Syntax-slot variables and non-stereo decode}
	\label{s:variables}

	For a fixed skeleton \(S\), the remaining choices are finite-domain variables.
	It is often convenient to package the non-directional choices into a base
	assignment \(u\), and then treat directional marker choices separately.

	\[
	\begin{array}{>{\raggedright\arraybackslash}p{0.23\linewidth}
			>{\raggedright\arraybackslash}p{0.26\linewidth}
			>{\raggedright\arraybackslash}p{0.38\linewidth}}
		\text{variable} & \text{domain} & \text{meaning} \\
		\hline
		\(a_v\) & finite \(A_{v,S}\) & atom text for atom \(v\) \\
		\(\tau_v\) & \(\{\varnothing,@,@@\}\), restricted by \(P\) & tetrahedral token, written separately from \(a_v\) only for clarity \\
		\(b_s\) & finite \(B_{s,S}\) & non-directional bond text at tree or ring slot \(s\) \\
		\(\delta_k\) & \(Z=\{0,+1,-1\}\) & directional marker value at carrier slot \(k\) \\
		\(\lambda_e\) & finite \(\Lambda\) & normalized ring label for cotree edge \(e\)
	\end{array}
	\]

	The value \(\delta_k=0\) is the absence of a slash/backslash marker at an
	existing carrier slot.  It is useful to keep this absence in the same finite
	domain as the two marker orientations.  A slot that does not exist for a given
	skeleton is different: it is outside the active carrier set \(K_S\), or is made
	inactive by an activation variable.  Thus there are two notions of absence: no
	syntax slot exists, and a slot exists but is emitted without a directional
	marker.

	For atom text, the local table may be written
	\[
	a_v\in A_{v,S}
	\quad\Longrightarrow\quad
	\Decode_{\rm atom}(v,S,a_v)=\alpha(v).
	\]
	For tree bond slots,
	\[
	b_s\in B_{s,S}
	\quad\Longrightarrow\quad
	\Decode_{\rm bond}(s,S,b_s,\delta_s)=\beta(e_s),
	\]
	where \(e_s\) is the graph edge represented by slot \(s\), ignoring the
	stereochemical meaning of \(\delta_s\) except insofar as the slash or backslash
	is syntactically a bond character.  For an omitted bond, the decode table is a
	relation involving the endpoint atom attributes and the edge attributes; this is
	where fixed aromaticity and fixed implicit-hydrogen conventions are used.

	For ring edges it is safer to use a two-endpoint relation, because a policy may
	allow a bond symbol at either endpoint or at both endpoints.  If
	\(e=\{u,v\}\in E_{\Ring}\) has endpoint slots \(s=(e,u)\) and \(t=(e,v)\), impose
	\[
	\mathsf{RingBond}_e(S,b_s,\delta_s,b_t,\delta_t)=\beta(e).
	\]
	If the policy fixes a unique convention, this relation is a singleton.  If the
	parser's non-stereo decode is not purely local, replace the preceding local
	tables by one finite relation
	\[
	\mathsf{NonStereoDecode}_G(S,u),
	\]
	which asserts that the rendered non-stereo syntax decodes to \((V,E,\alpha,\beta)\).
	Since \(G\), \(S\), and all text domains are finite, this global relation is
	still finite.  The local-table presentation is a convenient special case, not a
	hidden assumption.

	\section{Ring-label constraints}
	\label{s:ringlabels}

	For each cotree edge \(e\in E_{\Ring}\), the two endpoint occurrences receive a
	common label \(\lambda_e\in\Lambda\).  Let \(q_e^-\) and \(q_e^+\) be the first
	and second of the two endpoint occurrences in lexical order.  Reuse of a label
	is legal only after the previous use has closed.  Thus for distinct ring edges
	\(e\) and \(f\), impose
	\[
	\lambda_e=\lambda_f
	\quad\Longrightarrow\quad
	q_e^+<q_f^-\quad\text{or}\quad q_f^+<q_e^- .
	\]
	Equivalently, same-label intervals may not overlap.  This is the essential
	bounded reusable-label condition: a parser state never contains two open ring
	closures with the same label, and a closing occurrence is paired with the
	currently open occurrence of that label.

	There are also endpoint-specific correctness conditions, all finite:
	\begin{enumerate}[label=\textup{(R\arabic*)},leftmargin=2.5em]
		\item each ring edge has exactly two endpoint events, incident with the two
		endpoints of that edge;
		\item the two endpoint events use the same label text, and no tree edge uses a
		ring label;
		\item the bond text at the first and second endpoint is compatible with
		\(\beta(e)\), as enforced by \(\mathsf{RingBond}_e\);
		\item if \(P\) normalizes labels, for example by requiring the least currently
		free label, the normalization is imposed as the finite order constraint
		\[
		\lambda_e
		=\min\{\lambda\in\Lambda:
		\lambda\text{ is inactive immediately before }q_e^-\} .
		\]
	\end{enumerate}
	No additional procedural pairing rule is needed in the bounded model.  The usual
	parser stack or table for ring labels is represented declaratively by these
	finite interval and endpoint constraints.

	\section{Potential stereo sites}
	\label{s:potential-sites}

	The cleanest definition of a potential stereo site is parser-defined and
	policy-relative.  A graph-theoretic description is useful as a candidate
	generator, but it is not the definition; and a presentation policy is used to
	restrict which syntactic carriers are allowed, but it does not by itself define
	what the parser would regard as a stereo site.

	For a skeleton \(S\), let \(Y_S\) be the finite set of all syntax-slot
	assignments satisfying the non-stereo graph constraints, before imposing stereo
	preservation or annotation maximality.  The declared parser semantics provides a
	finite set of stereo-output coordinates
	\[
	\Site_P(G,S)
	\]
	and, for each \(q\in\Site_P(G,S)\), a finite value set
	\[
	\St_q=\{\None\}\cup \St_q^\times .
	\]
	The value \(\None\) means that no stereo is assigned at coordinate \(q\), while
	\(\St_q^\times\) contains the possible nontrivial stereo states.  Examples are
	the two tetrahedral states at a tetrahedral atom and the two directional states
	at an eligible double-bond site.

	\begin{definition}[Potential stereo site]
		For the bounded parser semantics and policy \(P\), a coordinate
		\(q\in\Site_P(G,S)\) is a potential stereo site for skeleton \(S\) if
		\[
		\exists y\in Y_S
		\quad \text{such that}\quad
		\parse_P(\render(S,y))\text{ is defined and }
		\parse_{P,q}(\render(S,y))\in\St_q^\times .
		\]
		The set of such sites is denoted \(Q_P(G,S)\).  In the single CSP over all
		skeletons, one uses the finite union of these sets together with activation
		constraints saying which sites and carrier scopes are present for the chosen
		skeleton.
	\end{definition}

	This definition has two advantages.  First, no-accidental-stereo constraints are
	placed exactly where the parser could assign stereo, not merely where a chemist
	could draw a graph-theoretic stereocentre.  Second, the definition automatically
	adapts to syntax policy.  If a carrier is forbidden by the policy, it cannot
	make a site potential.  If a marker is permitted at a ring endpoint, the local
	orientation and ring endpoint position are part of the relevant parser relation.

	For tetrahedral stereo, potential sites outside \(T\) are normally suppressed by
	forcing \(\tau_v=\varnothing\).  For directional stereo, potential sites are
	handled by finite relations over slash/backslash carrier slots, as described
	next.

	\section{Tetrahedral stereo constraints}
	\label{s:tetrahedral}

	For a tetrahedral centre \(v\in T\), let \(L_v\) be its set of four ligand
	occurrences.  A ligand occurrence may be a neighbouring graph atom or a fixed
	implicit-hydrogen occurrence.  The skeleton and ring endpoint order induce the
	local SMILES ligand order
	\[
	\theta_v(S)\in\operatorname{Seq}(L_v),
	\]
	where \(\operatorname{Seq}(L_v)\) is the finite set of length-four orderings of
	\(L_v\).  The order includes the incoming parent atom when \(v\) is not a root,
	the ring closure positions and child positions following the atom token, and the
	policy's convention for implicit hydrogens.  Ring digits matter here through
	their lexical position on the chiral atom, not through the later endpoint to
	which the digit is paired.

	For each specified centre \(v\in T\), define a finite relation
	\[
	\Tet_v
	\subseteq
	\operatorname{Seq}(L_v)\times \{@,@@\}\times\{\chi_v\} .
	\]
	The tetrahedral constraint is
	\[
	(\theta_v(S),\tau_v,\chi_v)\in\Tet_v .
	\]
	Equivalently, after choosing a reference ligand order \(\rho_v\), the same
	constraint can be written as a parity equation
	\[
	\operatorname{parity}(\theta_v(S),\rho_v)
	\cdot \operatorname{sign}(\tau_v)
	=\operatorname{sign}(\chi_v),
	\]
	with all convention choices absorbed into \(\Tet_v\).  For atoms not carrying
	specified tetrahedral stereo, impose \(\tau_v=\varnothing\), or more generally
	impose the parser-defined value \(\None\) at the corresponding potential site.

	Thus changing the root, branch order, continuation child, or ring digit position
	may change whether \(@\) or \(@@\) is required, but the change is captured by a
	finite local table.

	\section{Directional stereo and no accidental stereo}
	\label{s:directional}

	For a fixed skeleton \(S\), let \(K_S\) be the finite set of active syntax slots
	that are eligible to carry a directional marker.  Each \(k\in K_S\) has a
	variable
	\[
	\delta_k\in Z:=\{0,+1,-1\}.
	\]
	The actual marker characters \(\mathtt{/}\) and \(\backslash\) are obtained from \(\delta_k\) by
	a fixed finite rendering table that also knows the local text direction of the
	slot.  The sign is not an absolute geometric meaning of the printed character;
	it is a normalized syntactic value relative to the chosen slot orientation.

	For each potential directional site \(q\in Q_P(G,S)\), let
	\(K_{q,S}\subseteq K_S\) be the carrier slots whose markers can affect the parsed
	stereo value at \(q\).  The declared parser semantics gives a finite relation
	\[
	\Dir_{q,S}
	\subseteq
	Z^{K_{q,S}}\times \St_q .
	\]
	Assignments outside this relation are syntactically invalid, contradictory, or
	outside the bounded semantics.  If \(q\) is specified in the input with value
	\(\omega_q\), impose
	\[
	(\delta|_{K_{q,S}},\omega_q)\in\Dir_{q,S} .
	\]
	If \(q\) is a potential site but is not specified in the input, impose the
	no-accidental-stereo constraint
	\[
	(\delta|_{K_{q,S}},\None)\in\Dir_{q,S} .
	\]
	This is the clean form of ``do not introduce extra stereo''.  It is parser-defined:
	it asks that the declared parser assign no stereo at that output coordinate.

	No-accidental stereo is local in the following precise sense.  For ordinary
	alkene, imine, and oxime-style directional stereo, each constraint has a finite
	scope consisting of the carrier slots around one potential site.  However, the
	whole problem is not independent by site, because the same carrier slot may be
	adjacent to more than one potential site, especially in conjugated or
	shared-carrier situations.  The local relations form overlapping hyperedges in a
	global CSP.  If a particular bounded parser semantics performs conflict
	resolution or cleanup over a larger directional component, replace the one-site
	relations by a component relation
	\[
	\Dir_{C,S}
	\subseteq
	Z^{K_{C,S}}\times \prod_{q\in C}\St_q,
	\]
	where \(C\) is a finite set of coupled sites.  Then no-accidental stereo is the
	coordinate constraint that unspecified \(q\in C\) receive value \(\None\).  Thus
	larger-than-local parser behaviour does not obstruct the finite model; it merely
	enlarges the finite relation's scope.

	\subsection{Parity form and orientation factors}
	\label{s:parity}

	For the ordinary pairwise cases, \(\Dir_{q,S}\) is often expressible as a signed
	parity CSP.  Suppose a potential site \(q\) is represented by a double bond with
	endpoints \(u\) and \(v\), and carrier sets \(K_{q,u}\) and \(K_{q,v}\).  For
	each carrier \(k\in K_{q,S}\), let
	\[
	\kappa_{q,k}\in\{+1,-1\}
	\]
	be the orientation factor converting the local direction of the printed slot
	into the orientation seen from the relevant endpoint of the double bond.  For
	\(\delta_k\ne0\), set
	\[
	\sigma_{q,k}:=\kappa_{q,k}\delta_k .
	\]
	Then the nontrivial state at \(q\) can be enforced by finite constraints such as
	\begin{align*}
		&\exists k\in K_{q,u},\;\exists \ell\in K_{q,v}:
		\delta_k\ne0,
		\;\delta_\ell\ne0,\\
		&\sigma_{q,k}\sigma_{q,\ell}=\eta_{q,k,\ell},
		&&k\in K_{q,u},\ \ell\in K_{q,v},\
		\delta_k\delta_\ell\ne0,\\
		&\sigma_{q,k}\sigma_{q,k'}=\eta^u_{q,k,k'},
		&&k,k'\in K_{q,u},\
		\delta_k\delta_{k'}\ne0,\\
		&\sigma_{q,\ell}\sigma_{q,\ell'}=\eta^v_{q,\ell,\ell'},
		&&\ell,\ell'\in K_{q,v},\
		\delta_\ell\delta_{\ell'}\ne0 .
	\end{align*}
	The constants \(\eta\) are determined by the input stereo descriptor and the
	chosen convention.  This parity representation is sufficient for ordinary
	pairwise directional stereo once the parser semantics really is only assigning
	relative side information from carrier signs.  The primitive relation
	\(\Dir_{q,S}\) is nevertheless the safer mathematical object: it can also encode
	endpoint degeneracies, ring-endpoint orientation, missing-carrier rules,
	contradictory marker patterns, and any bounded parser-specific interpretation
	that remains within the declared dialect.

	A common pitfall is to treat \(\mathtt{/}\) or \(\backslash\) as carrying an absolute
	meaning independent of traversal.  It does not.  The same printed token may
	contribute different normalized signs depending on whether the syntax slot is
	written toward or away from the double-bond endpoint, whether it occurs at a
	ring-closure endpoint, and which adjacent potential site is being considered.
	The orientation factors \(\kappa_{q,k}\) are therefore part of the semantics, not
	an implementation detail.

	\section{Annotation policies}
	\label{s:annotation-policy}

	The most general and cleanest formal type of an annotation policy is a predicate
	on complete assignments:
	\[
	\Ann_P(S,u,\delta)\in\Bool,
	\]
	where \(S\) is the skeleton, \(u\) is the non-directional base assignment, and
	\(\delta\) is the directional marker assignment on the active carrier set
	\(K_S\).  This predicate can include purely presentational rules, semantic
	maximality rules, and normalization rules.  Treating it as a predicate avoids a
	false distinction between ``generation'' and ``validation'': it is simply one
	block of the finite constraint instance.

	Three common presentations are extensionally equivalent only in special cases.
	A pre-solver generator of allowed slots is adequate for local eligibility rules,
	such as forbidding markers on non-single bonds or outside a selected carrier
	set.  It is not adequate for policies whose truth depends on whether the rest of
	the stereo constraints are jointly satisfiable.  An optimization or maximality
	layer is the honest formulation for ``as many markers as possible'' or
	``inclusion-maximal marker support''.  A hard predicate on complete assignments
	is the common mathematical abstraction that covers both.

	For fixed \((S,u)\), let
	\[
	\OK_{S,u}(\delta)
	\]
	be the conjunction of the directional preservation constraints and the
	no-accidental-stereo constraints.  Let
	\[
	M(\delta):=\{k\in K_S:\delta_k\ne0\}
	\]
	be the marker support.  The support-wise maximal policy is
	\[
	\begin{split}
		\Max^{\subseteq}_{S,u}(\delta)
		:\Longleftrightarrow{}&
		\OK_{S,u}(\delta)\quad\text{and}\quad\\
		&\nexists\delta'\in Z^{K_S}
		\text{ such that }
		\OK_{S,u}(\delta')
		\text{ and }
		M(\delta)\subsetneq M(\delta').
	\end{split}
	\]
	This is the natural reading of ``emit every eligible carrier marker that can be
	included in a surviving assignment'', where survival means satisfying all stereo
	and no-extra-stereo constraints for the same base rendering.  It may admit
	several incomparable maximal supports.

	If the intended policy is ``maximize the number of markers'', use instead
	\[
	\Max^{\#}_{S,u}(\delta)
	:\Longleftrightarrow
	\OK_{S,u}(\delta)
	\text{ and }
	|M(\delta)|=
	\max\{|M(\delta')|:\OK_{S,u}(\delta')\}.
	\]
	If the intended policy is writer-like uniqueness, add a tie-breaker, for example
	lexicographic selection among the cardinality-maximal supports.  These variants
	are not equivalent in general.  Support-wise maximality preserves all
	inclusion-maximal marker patterns; cardinality maximality discards smaller
	inclusion-maximal patterns; canonicalization keeps one selected pattern.

	Because \(K_S\) is finite, all of these policies are finite.  The compilation
	cost differs.  Cardinality maximality is naturally a MaxSAT or optimization
	layer.  Inclusion maximality is naturally a maximal-model or QBF-style
	condition.  It can be compiled into ordinary SAT by finite expansion or by
	adding blocking constraints during enumeration, but in general one should not
	pretend that there is a uniformly small CNF encoding unless additional monotonic
	structure is proved.  The optimization/maximality layer is the honest formal
	description.

	\section{The finite constraint instance and support}
	\label{s:instance}

	Let \(\calC(G,P)\) be the finite constraint instance whose variables are the
	skeleton variables and all syntax-slot variables, and whose constraints are:
	\begin{enumerate}[label=\textup{(C\arabic*)},leftmargin=2.5em]
		\item traversal-skeleton constraints;
		\item atom, bond, incidence, component, and edge-coverage constraints;
		\item ring-label endpoint, interval, and optional normalization constraints;
		\item tetrahedral parser relations;
		\item directional parser relations, including no-accidental-stereo constraints;
		\item the annotation-policy predicate \(\Ann_P\);
		\item optional duplicate-witness constraints if the enumerator must output each
		string exactly once.
	\end{enumerate}
	All variables range over finite domains.  Therefore \(\calC(G,P)\) is a finite
	CSP, or after Boolean encoding, a finite SAT/MaxSAT/QBF-style object depending
	on how the annotation policy is represented.

	Let \(\Sol\calC(G,P)\) be the satisfying assignments.  The renderer gives a map
	\[
	\render:\Sol\calC(G,P)\longrightarrow\Sigma_P^*,
	\]
	where \(\Sigma_P\) is the finite token alphabet of the bounded language.  The
	exact support is
	\[
	L(G,P):=\render(\Sol\calC(G,P)).
	\]
	This image-of-render definition is mathematically clean.  It defines support as
	a set of strings, not as a set of derivations or labelled graph witnesses.

	The map \(\render\) need not be injective.  Symmetric molecules may have distinct
	root choices or traversal witnesses that render the same string.  A parsed
	string is compared with the input graph only up to labelled stereochemical graph
	isomorphism, and a completeness proof may choose any one such isomorphism to
	read off a witness.  No quotient by graph automorphisms is needed in the
	definition of \(L(G,P)\); automorphisms only create duplicate witnesses in the
	preimage of the same string.  If duplicate-free enumeration is required, add a
	finite canonical-witness constraint, such as selecting the lexicographically
	least event-code witness among all assignments in the relevant automorphism
	orbit that render the same token string, or enumerate the image and deduplicate
	externally.  The former is declarative; the latter is an enumeration strategy.
	Neither is a post-hoc parse filter.

	\section{Exactness theorem}
	\label{s:theorem}

	\begin{theorem}[Exact support]
		Let \(G\) be a fixed finite molecular graph and let \(P\) be a finite bounded
		SMILES policy satisfying the assumptions above.  Let \(\calC(G,P)\) be the
		constraint instance just defined.  Then
		\[
		L(G,P)=\render(\Sol\calC(G,P))
		\]
		is exactly the set of strings in the declared sublanguage whose parse under the
		declared semantics is isomorphic to \(G\) as a labelled stereochemical graph and
		whose presentation satisfies the declared annotation policy.
	\end{theorem}

	\begin{proof}
		For soundness, let \(x\in\Sol\calC(G,P)\).  The traversal constraints give a
		well-formed depth-first fragment rendering with balanced branches, legal dots,
		legal atom events, and paired ring-closure endpoints.  The coverage and
		incidence constraints ensure that every vertex of \(G\) appears once, every edge
		of \(G\) appears once either as a tree bond or as a paired ring closure, and no
		non-edge is emitted.  The atom and bond decode relations ensure that the parsed
		non-stereochemical labels are \(\alpha\) and \(\beta\).  The ring-label interval
		constraints ensure that any reused label is inactive before being reused, and
		therefore paired with the intended cotree edge under the declared parser
		semantics.

		For each tetrahedral centre, the local SMILES ligand order induced by the
		rendered syntax is \(\theta_v(S)\), and the relation \(\Tet_v\) forces the parsed
		tetrahedral state to be \(\chi_v\).  For each specified directional site, the
		relation \(\Dir_{q,S}\) forces the parsed directional state to be the input value
		\(\omega_q\).  For each unspecified potential site, the same parser relation
		forces the parsed value to be \(\None\), so no additional directional stereo is
		introduced.  Finally, \(\Ann_P\) keeps only the presentations allowed by the
		annotation policy.  Thus \(\render(x)\) is syntactically valid and parses to
		\(G\) with exactly the prescribed stereochemistry.

		For completeness, let \(w\) be a string in the declared sublanguage whose parse
		under \(P\) is isomorphic to \(G\) and whose presentation satisfies the policy.
		Choose one parse isomorphism from the parsed molecule to \(G\).  Reading \(w\)
		from left to right gives the roots, parent map, branch and continuation choices,
		ring endpoint positions, ring labels, atom texts, bond texts, tetrahedral tokens,
		and directional marker values.  Since \(w\) belongs to the declared sublanguage,
		these values lie in the finite domains of \(P\).  Since \(w\) parses to \(G\), the
		coverage, decode, tetrahedral, directional, and no-accidental-stereo relations
		are satisfied.  Since \(w\) obeys the annotation policy, \(\Ann_P\) is satisfied.
		The read-off data therefore form a satisfying assignment \(x\) with
		\(\render(x)=w\).
	\end{proof}

	The proof uses no hidden parser determinism beyond the declared semantics.  If
	the parser semantics is a relation rather than a function, the same proof works
	after replacing ``the parse'' by the selected relational criterion, for example
	``there exists a parse isomorphic to \(G\)'' or ``all legal parses are isomorphic
	to \(G\)''.  The criterion must be chosen explicitly.

	\section{Equivalent formalisms}
	\label{s:formalism}

	The formulation above can be expressed in several standard languages.
	\begin{enumerate}[label=\textup{(F\arabic*)},leftmargin=2.5em]
		\item An attributed graph-traversal grammar plus a finite CSP is the most direct
		specification.  The grammar supplies the legal DFS event skeletons; the CSP
		supplies finite semantic relations.
		\item A regular tree grammar with constraints can generate the ordered traversal
		trees, with equality and finite-table constraints for ring labels, slot text,
		and stereo attributes.
		\item A constraint automaton or finite-state transducer with registers can be
		used after bounding the graph, labels, and syntax domains.  Full compilation
		produces an ordinary finite automaton because the language is finite.
		\item Monadic second-order logic on a finite ordered structure can describe
		spanning forests, event orders, matching ring endpoints, and finite label
		relations.  This is elegant for meta-theorems but less transparent as a working
		SMILES specification.
		\item A graph-to-string transducer with finite attributes gives another
		presentation of the same object.
	\end{enumerate}
	These formalisms are interchangeable at the level of existence.  The attributed
	traversal grammar plus CSP is preferable because it keeps the graph traversal,
	syntax slots, and parser semantics visibly separated.

	\section{Complexity and enumeration target}
	\label{s:complexity}

	The finite model is an exact specification, not a claim of polynomial-time
	enumeration.  Complexity statements should distinguish three quantities:
	\[
	\begin{aligned}
		N_{\rm wit}&=|\Sol\calC(G,P)|,\\
		N_{\rm str}&=|L(G,P)|,\\
		N_{\rm can}&=\text{number of canonical witnesses, if imposed}.
	\end{aligned}
	\]
	Enumerating satisfying assignments is naturally output-sensitive in
	\(N_{\rm wit}\), but this may be the wrong measure if many witnesses render the
	same string.  Exact support enumeration as a language problem should be measured
	output-sensitively in \(N_{\rm str}\).  To obtain such an enumeration one either
	enumerates canonical witnesses or enumerates assignments while maintaining a set
	of already rendered strings.  The former pushes duplicate elimination into the
	constraint model; the latter performs duplicate elimination as an enumeration
	bookkeeping step.  Neither changes the definition of the support.

	If the annotation policy is a plain finite predicate, Boolean encoding gives an
	ordinary SAT/CSP enumeration problem.  Cardinality-maximal policies are naturally
	MaxSAT or optimization problems.  Inclusion-maximal policies are naturally
	maximal-model or quantified conditions.  They remain finite, but compiling them
	into ordinary SAT without an optimization or oracle layer may require an
	exponential-size expansion in the number of carriers unless additional structure
	is available.

	\section{Semantic constraints and presentation constraints}
	\label{s:semantic-presentation}

	All constraints can be put into one finite conjunction, but it is conceptually
	useful to keep two roles separate.  Semantic constraints define the parser
	preimage of the labelled stereochemical graph: graph coverage, atom and bond
	decode, tetrahedral parity, directional decode, and no-accidental stereo.
	Presentation-policy constraints choose which strings inside that semantic
	preimage belong to the support: branch-order restrictions, fragment-order rules,
	ring-label normalization, atom-text style, and annotation maximality or
	canonicalization.

	This distinction is not required by the mathematics, but it is useful for
	proofs.  Soundness of the semantic block says that every rendered string parses
	to \(G\).  Soundness of the policy block says that every rendered string is in
	the chosen sublanguage.  Combining the two as one CSP is harmless as long as the
	target support is stated explicitly.

	\section{Where obstructions would arise}
	\label{s:obstructions}

	Under the bounded assumptions there is no fundamental obstruction.  The
	apparent hard cases are finite relations or finite optimality predicates:
	\[
	\begin{array}{>{\raggedright\arraybackslash}p{0.38\linewidth}
			>{\raggedright\arraybackslash}p{0.49\linewidth}}
		\text{issue} & \text{finite declarative treatment} \\
		\hline
		branch order, root choice, or ring digit position changes \(@\)/\(@@\) & local
		ligand-order table \(\Tet_v\) \\
		slash/backslash token has traversal-dependent orientation & slot- and
		site-specific orientation factor \(\kappa_{q,k}\) or finite relation
		\(\Dir_{q,S}\) \\
		markers interact through conjugated or shared-carrier systems & overlapping
		finite directional constraints, or one component relation \(\Dir_{C,S}\) \\
		markers may accidentally assign unspecified stereo & parser-defined potential
		sites constrained to value \(\None\) \\
		ring labels are reused & same-label active-interval non-overlap plus endpoint
		and bond-text compatibility \\
		symmetric graphs create duplicate witnesses & define support as
		\(\im(\render)\), or add canonical-witness constraints \\
		maximal annotation is ambiguous & choose inclusion-maximal, cardinality-maximal,
		or canonicalized policy explicitly
	\end{array}
	\]

	Obstructions appear when finiteness or explicit semantics are dropped:
	\begin{enumerate}[label=\textup{(O\arabic*)},leftmargin=2.5em]
		\item unbounded ring-label syntax, atom-text syntax, extension fields, or
		arbitrary redundant spellings can make the support infinite;
		\item aromaticity, valence, implicit hydrogens, or stereocentre cleanup performed
		by an external parser must be encoded as finite relations or used as oracles;
		\item a target defined as ``whatever a particular writer emits after repair'' is
		a writer-specific target unless the repair rules are formalized;
		\item features outside the bounded dialect, such as enhanced stereo groups,
		non-tetrahedral stereo, atropisomerism, or unknown/relative stereo, require their
		own finite semantics if admitted;
		\item an informal annotation policy does not define a support.
	\end{enumerate}

	\section{Weakest assumptions for avoiding a post-hoc parse filter}
	\label{s:weakest-assumptions}

	The theorem does not require the exact RDKit writer, nor does it require a
	canonical SMILES algorithm.  The weakest useful assumptions are these.
	\begin{enumerate}[label=\textup{(A\arabic*)},leftmargin=2.5em]
		\item The input object \(G\) is finite.
		\item The syntax policy has finite domains for every atom, bond, ring-label,
		fragment, branch, traversal, and stereo-annotation choice.
		\item The traversal grammar emits only syntactically well-formed event skeletons,
		including balanced branches and well-formed ring endpoint events.
		\item The parser semantics relevant to the bounded dialect is explicitly given
		as a finite relation or partial function on the event and slot variables.
		\item The desired support is exactly the set of rendered strings satisfying that
		finite semantic relation and the finite presentation policy.
	\end{enumerate}
	Under these assumptions, the parse-equivalence predicate is already part of the
	finite constraint instance.  Running a parser afterwards may be useful for
	software testing, but it is not part of the mathematical enumeration procedure
	and is not needed to reject invalid outputs.

	\section{Conclusion}
	\label{s:conclusion}

	For a fixed finite stereochemical graph and a bounded RDKit/OpenSMILES-like
	SMILES dialect with explicit syntax and parser semantics, exact support
	enumeration is a single finite declarative problem.  The preferred model is an
	attributed ordered graph-traversal grammar coupled to finite-domain constraints
	over syntax slots.  Potential stereo sites are defined by the parser semantics
	relative to the policy.  No-accidental stereo is enforced by finite relations
	whose scopes may overlap globally through shared carrier variables.  Annotation
	policies are best treated as predicates on complete assignments, with
	maximality or optimization layers used when the policy itself is maximal rather
	than purely local.

	The support is cleanly defined as the image of the rendering map on satisfying
	assignments.  Canonical witnesses, automorphism quotients, and duplicate
	elimination are enumeration refinements, not prerequisites for the exactness
	theorem.  Once all parser and policy choices are finite and explicit, no
	post-hoc parse filtering, writer-specific repair rule, or feature-specific
	mini-oracle is mathematically forced.

\end{document}
